\section{Methods}

\subsection{Input class}

Fix constants $c_\kappa,c_\eta,B_{\rm in}>0$. An input consists of a nonsingular
$M\in\Q^{n\times n}$ with $n\ge1$, a right-hand side $b\in\Q^n$ and a tolerance $\eta$, with
\begin{equation}
 m=\nnz(M),\qquad \kappa_2(M)\le m^{c_\kappa},\qquad m^{-c_\eta}\le\eta\le1.
 \label{eq:methods-input}
\end{equation}
Every numerator and denominator has at most $B_{\rm in}\log_2(2+m)$ bits, and
indices have $O(\log m)$ bits. Nonsingularity implies $n\le m$. The output is a
rational vector $x$ with $\norm{Mx-b}_2\le\eta\norm b_2$. The encoder may
replace $\eta$ by $\eta_0=m^{-\lceil c_\eta\rceil}$, so that the tolerance never has to be
read to more than logarithmic precision. The empty system is solved directly.
Constants in all asymptotic bounds depend only on $c_\kappa$, $c_\eta$ and $B_{\rm in}$.

\subsection{Continuum realization}

The domain $\Omega$ is a rectangular box with integer side lengths, and the
field satisfies \cref{eq:maxwell} with $1\le\varepsilon(r)\le25$. Let
$\mathcal C$ denote the self-adjoint curl-curl operator defined by the closed
form $\langle\nabla\times E,\nabla\times F\rangle$ on
$H_0(\operatorname{curl};\Omega)$, let $X$ be its domain with the graph norm
$\norm E_X^2=\norm E_{L^2}^2+\norm{\mathcal CE}_{L^2}^2$, and set
$A=\mathcal C-\varepsilon$. For fixed $c_\kappa,c_\eta,B_{\rm in}$, there are a randomized
encoder, a decoder and explicit constants $C_{\rm res},c_A,c_\delta$ such that, with
$\ell=1+\lceil\log_2(2+m)\rceil$, every run of the encoder satisfies
\begin{align}
 |\Omega|&\le C_{\rm res}m^{3/2}\ell^{24},&
 N_{\rm desc}&\le C_{\rm res}m^{3/2}\ell^{64},\notag\\
 T_{\rm enc}+T_{\rm dec}&\le C_{\rm res}m^{3/2}\ell^{4096},&
 B_{\rm out}&\le C_{\rm res}\ell,\qquad \delta\ge m^{-c_\delta}.
 \label{eq:methods-continuum-cost}
\end{align}
Here $N_{\rm desc}$ counts the unit cells of the box together with all
coefficients and indices of the description, and $B_{\rm out}$ bounds the bit
length of every output coefficient. For $m\ge2$, with probability at least
$1-3/(64m)$, the operator $A:X\to L^2$ is invertible with
\begin{equation}
 \norm{A^{-1}}_{L^2\to L^2}\le L_b,\qquad
 \norm A_{X\to L^2}\norm{A^{-1}}_{L^2\to X}\le702L_b\le m^{c_A},
 \label{eq:methods-inverse}
\end{equation}
and every $E\in X$ with $\norm{AE-f}_{L^2}\le\delta\norm f_{L^2}$ decodes to a
valid answer for every allowed error in the moment messages, whose range and
precision both have $O(\ell)$ bits. The condition bound refers to the graph
norm, since $A$ is unbounded on $L^2$. The complete statement is Theorem~\ref{c:thm:main}
of the \SI, proved in Part~II.

\subsection{Yee realization}

The discrete target uses the Yee curl incidence matrix $C_{\rm Y}$ with unit
grid spacing, frequency $\omega=1/8$ and PEC termination; $A$ now denotes the
resulting Yee matrix. Every edge mass is
the average of the permittivities of its four adjacent cells, and
\begin{equation}
 A=C_{\rm Y}^{\mathsf T}C_{\rm Y}-\omega^2D_\varepsilon,
 \qquad 1\le\varepsilon_c\le25.
 \label{eq:methods-discrete}
\end{equation}
The grid is uniform, without coordinate stretching or absorbing layers, and
resolves the shortest wavelength in the material with ten cells. The
realization has $\Ot(m^{3/2})$ unknowns, the same encoding and decoding cost,
$O(\log m)$-bit coefficients and, with probability at least $1-3/(64m)$, a
full matrix with $\kappa_2(A)\le m^{c_A}$ for which every vector meeting
the residual tolerance decodes correctly. The decoder reads prescribed field
coordinates to $O(\log m)$ bits. The spectral analysis of the discrete model is
independent of the continuum one and does not rely on the resolution to
approximate the continuum equations (Theorem~\ref{d:thm:main} and Part~III
of the \SI). Its spectral input is a computer-assisted certificate for a
periodic Yee background. The finite inequalities of this certificate have been
checked in full; the final step, which binds the checked data to the exact Yee
operator near the band minimum, is being completed as part of an ongoing formal
verification, and the Yee result is conditional on it
(Section~\ref{si:sec:yee-certificate}).

\subsection{Logical network and routing}

The compiler normalizes and rounds the coefficients within the error budget,
expresses their arithmetic through a fixed alphabet of local relations, and
adds variables to bound the degree of every site by three. The original system
is embedded in a symmetric system, and sign mediators make all couplings of the
realized network positive, with the required signs appearing after
elimination of the mediators. Balanced routing then replaces every edge by a
nearest-neighbor path in a cubic grid of side $O(\sqrt{m\log m})$. Every step
preserves a polynomial inverse bound and a linear map from the residual of the
network to the residual of $Mx=b$ (\SI, Part~I).

\subsection{Band minimum, interactions and leakage}

The background of \cref{eq:background} has a simple Bloch branch with
$\lambda(k_*+q)=E_c+q^{\mathsf T}Hq+O(|q|^4)$ and $H$ positive definite.
Exact rational arithmetic identifies the branch in the finite corner cluster,
and estimates over all other Fourier components establish the expansion for
the continuum operator. An integer placement matrix adapts the routing grid to
the anisotropy of $H$.

Defects are smooth bumps of radius $L$ and strength of order $L^{-2}$, analyzed
through an auxiliary resolvent localized in momentum near $k_*$. Along the
three routing directions, the coupling between neighboring defects at
separation $d$ satisfies $h_j(d)\ge2^{-1024}L^2 8^{-d}$, while the full coupling
block between defects displaced by $\nu\ne0$ lattice steps satisfies
$\norm{G_h(\nu)}\le C_*L^5(3^{-d\norm \nu_1}+e^{-c_*L})$. The margin
$(1/3)^2<1/8$ lets the useful couplings dominate the summed unwanted ones after
diagonal tuning. The momentum cutoff leaves a residual in the actual Maxwell
equation bounded by $e^{-c_*L^{1/3}}$ times explicit polynomial factors.
Choosing $L=O(\ell^6)$, $d=O(\ell)$ and a boundary buffer of width $O(\ell^7)$
makes all errors smaller than the required inverse-polynomial budgets. The
uniform permittivity shift is sampled after all profiles are fixed, and a
count of the eigenvalues of the finite-box operator gives
\cref{eq:methods-inverse} (\SI, Part~II).

\subsection{Finite description and recovery}

The permittivity consists of the trigonometric background, dyadic shifts and
polynomial profiles on balls of radius growing with $m$. Calibration uses
finite algebra at $k_*$ and controlled Fourier tails and does not solve the
assembled problem. Sources and weighting functions use tensor polynomials of
coordinate degree $O(\ell^8)$, with $O(\ell)$ bits per coefficient and error
bounds for the complete functions. If $E_{\rm ref}$ is the reference field
constructed in the proof, with $\norm{AE_{\rm ref}-f}_{L^2}\le r_{\rm ref}$, then
every field meeting the residual tolerance satisfies
\begin{equation}
 \norm{E-E_{\rm ref}}_{L^2}\le L_b\bigl(\delta\norm f_{L^2}+r_{\rm ref}\bigr).
 \label{eq:methods-comparison}
\end{equation}
The decoder reads the moments $z_i=v_i\langle E,\widehat a_i\rangle$, whose
real test fields $\widehat a_i$, weights $v_i$ and tolerances are part of the
output. \Cref{eq:methods-comparison}, the approximation bounds for the test
fields and the residual map of the logical network give \cref{eq:recovery}.

\subsection{Algorithm transfer and conditional limits}

For the continuum transfer, $N$ counts the volume and the description, and $B_{\rm out}$
the coefficient and message lengths. The solver is assumed to satisfy
\begin{equation}
 T_{\rm Max}(N,B_{\rm out},\kappa_X,\delta)
 =O\!\left(N^\theta B_{\rm out}\,\polylog(N,\kappa_X,\delta^{-1})\right),
 \label{eq:methods-solver}
\end{equation}
where $\kappa_X$ is the condition bound of \cref{eq:methods-inverse}, and to
return the prescribed moments of a residual field with a fixed positive
probability, the cost of computing the moments included. A polynomial
dependence on $\kappa_X$, as in unpreconditioned iterative methods, does not
meet this assumption. Solver calls are capped, candidate answers are checked
against $Mx=b$, and repetition amplifies the success probability, which gives
\cref{eq:transfer} in the bit model.

The conditional statement uses a real-number formulation: a full-column-rank
$F_0$ with $m$ nonzeros and $\kappa_2(F_0)\le m^{c_\kappa}$, and the target
\begin{equation}
 \norm{F_0x-\Pi_{F_0}b}_2^2\le n^{-10c_\kappa}\norm{\Pi_{F_0}b}_2^2,
 \label{eq:methods-realram}
\end{equation}
where $\Pi_{F_0}$ projects onto the column space. A reduction through a sparse
saddle-point system preserves this criterion even when
$\norm{\Pi_{F_0}b}\ll\norm b$. The hypothesis states that no randomized
algorithm solves every such instance with probability at least $2/3$ in time
$\Ot(m^{\gamma'})$ for any $\gamma'<\gamma$; it is a randomized
exponent-gap version of the hypothesis used for structured linear systems and
packing linear programs~\citep{kyngzhang2017,kyng2020}. For $\gamma>3/2$ it
excludes $\theta<2\gamma/3$ in \cref{eq:methods-solver}. Conjugate gradients in
exact arithmetic require $O(mn)\le O(m^2)$ operations, and their convergence
under the condition-number bound limits the hypothesis further, so that
\begin{equation}
 \tfrac32<\gamma\le\min\{2,1+c_\kappa\}.
 \label{eq:methods-gamma}
\end{equation}
Faster algorithms for the same model would restrict this range further
(Theorems~\ref{d:thm:algorithm-transfer} and~\ref{c:thm:transfer} and
Part~IV of the \SI).

\subsection{Computer-assisted steps}

The finite part of the continuum band analysis -- the decomposition of the
corner cluster, orthogonality, ordering of eigenvalues, velocity and Hessian
identities and placement inequalities -- is checked in exact rational
arithmetic. The discrete spectral analysis uses exact matrix inequalities on a
covering of the Brillouin zone. In both cases, the analytic arguments in the
\SI\ show how these finite inequalities imply the required spectral
properties (\SI, Part~V). The formal verification of the continuum
construction in Lean is substantially complete, with proofs for the continuum
operator, band and defect estimates, leakage, finite representations and
residual-to-solution recovery. A theorem map and verification records in the
companion repository specify the checked statements and their dependencies.
